\section{Prezentarea predictorului OGhel}

Predictia ramificatiilor reprezinta una dintre tehnicile fundamentale de optimizare a performantei procesoarelor moderne cu pipeline adanc.
In absenta unei predictii corecte, procesorul este nevoit sa goleasca pipeline-ul si sa reia executia de la adresa corecta,
operatie cunoscuta sub numele de \textit{branch misprediction penalty}. De-a lungul timpului, cercetatorii au propus o
varietate de strategii cu complexitati si acurateti diferite.

\subsection*{Predictori statici si bazati pe Markov.}
Cei mai simpli predictori sunt cei statici: \textit{always taken} si \textit{always not-taken}, care nu tin cont de istoricul executiei.
O treapta superioara o reprezinta predictorul bazat pe \textbf{lanturi Markov}, care modeleaza probabilitatea urmatorului rezultat
al unui salt in functie de cele $k$ rezultate anterioare ale aceluiasi salt. Desi intuitiv, modelul Markov este limitat de faptul
ca ia in considerare exclusiv istoricul local al unui singur salt, ignorand corelatiile cu alte ramificatii din program.

\subsection*{Predictori globali -- familia GAg/GAp/PAp.}
Schema \textbf{GAg} (\textit{Global history, All branches, global prediction table}) introduce notiunea de
\textit{registru de istorie globala} (GHR -- \textit{Global History Register}), in care se acumuleaza
rezultatele ultimelor $W$ salturi, indiferent de adresa acestora. Indicele in tabela de predictie se
formeaza direct din GHR. \textbf{GAp} pastreaza istoricul global, dar foloseste tabele de predictie
per-ramificatie, reducand aliasul dintre intrari. \textbf{PAp} merge mai departe, combinand un
registru de istorie per-ramificatie cu tabele de predictie individuale. Aceste scheme ofera o
acuratete ridicata pentru programe cu corelatii inter-ramificatii, insa tabela creste exponential
cu adancimea istoricului.

\subsection*{Predictorul GShare si predictori combinati.}
\textbf{GShare} aplica operatia XOR intre GHR si adresa PC a saltului pentru a indexa tabela de predictie,
reducand aliasul fara a creste semnificativ dimensiunea structurii. Predictorii \textbf{combinati} (de tip \textit{tournament}),
cum este cel implementat in SimpleSim sub optiunea \texttt{comb}, mentin doua predictori de baza si o \textit{meta-tabela}
care decide dinamic caruia sa ii acorde incredere pentru fiecare salt in parte.

\subsection*{Predictori bazati pe perceptron.}
O directie moderna si importanta este reprezentata de predictorii bazati pe \textbf{perceptron}. Fiecare ramificatie este asociata
unui vector de ponderi; rezultatul predictiei este determinat de semnul produsului scalar dintre vectorul de ponderi si vectorul de
istorie globala (codificat ca $\pm 1$). Antrenamentul consta in incrementarea sau decrementarea ponderilor corespunzatoare in functie
de rezultatul real. Avantajul principal este capacitatea de a captura corelatii liniare pe istorii foarte lungi, la un cost de stocare
rezonabil. Limita fundamentala a perceptronului simplu consta insa in incapacitatea de a invata functii nelineare: de exemplu, paritatea
unui sir de biti (corelatii XOR) nu poate fi reprezentata de un model liniar.

\subsection*{Predictorul O-GEHL.}
\textbf{O-GEHL} (\textit{O-Geometric History Length}) depaseste limitarile mentionate mai sus printr-o serie de mecanisme specifice.
In locul unei singure tabele indexate dupa un istoric de lungime fixa, O-GEHL utilizeaza \textbf{mai multe tabele independente} $T_0, T_1, \ldots, T_{M-1}$,
fiecare indexata dupa un istoric de lungime diferita. Lungimile istoricului formeaza o \textbf{progresie geometrica}: $L_i = \lfloor \alpha^i \cdot L_0 \rceil$,
unde $\alpha > 1$ este un factor de scala. Aceasta alegere asigura acoperirea uniforma a fenomenelor de corelatie la scari multiple ale
istoricului, de la corelatie pe termen scurt pana la corelatie pe termen lung. Predictia finala se obtine prin \textbf{suma algebrica}
a valorilor contoarelor saturate extrase din fiecare tabela -- un mecanism similar votului ponderat al unui ansamblu de perceptroni,
dar fara a fi limitat la corelatii liniare, datorita indexarii diferentiate. Aceasta suma este comparata cu zero: daca este pozitiva,
saltul este prezis ca efectuat (\textit{taken}), altfel ca neefectuat. Spre deosebire de predictori globali simpli, \textbf{aliasul}
este gestionat prin diversificarea functiilor de indexare: fiecare tabela foloseste o combinatie diferita de biti din PC si din GHR,
reducand probabilitatea ca doua ramificatii distincte sa partajeze acelasi contor in toate tabelele simultan. In concluzie, O-GEHL
reuneste avantajele istoricului global (capturarea corelatiilor inter-ramificatii), ale lungimilor geometrice (acoperire multi-scala),
si ale votului prin suma (robustete la zgomot si alias), oferind o acuratete superioara predictorilor clasici la un buget de hardware comparabil.

\subsection{Parametrii predictorului GEHL}

Proiectarea unui predictor GEHL implică patru grade de libertate principale. Alegerea valorilor optime pentru fiecare parametru constituie un compromis
între acuratețe, cost hardware și latență.

\subsubsection*{Numărul de tabele $M$}

GEHL utilizează $M$ tabele independente, fiecare indexat după un istoric de lungime
diferită. Tabelul $T_0$ nu folosește istoric global (\textit{bias table}, $L[0]=0$),
iar tabelele $T_1,\ldots,T_{M-1}$ utilizează lungimi crescătoare.
Această diversificare permite capturarea atât a corelațiilor pe termen scurt, cât și
a celor pe termen lung.

Studiile empirice arată că intervalul $M \in [4,\,12]$ oferă cel mai bun raport
acuratețe\,/\,cost. Sub 4 tabele granularitatea istoricului este insuficientă; peste 12
câștigul marginal de precizie nu justifică costul hardware suplimentar
(\textit{diminishing returns}).

\subsubsection*{Contoarele saturate pe $k$ biți}

Fiecare intrare dintr-un tabel stochează un contor saturat pe $k$ biți, care codifică
„încrederea" predictorului pentru un anumit pattern de istoric.
Valoarea implicită $k=4$ reprezintă un compromis bun: $k=2$ este prea instabil, iar
$k\geq8$ risipește spațiu fără câștig semnificativ de acuratețe.

\subsubsection*{Seria geometrică a lungimilor de istoric}

Lungimile de istoric urmează o \textbf{progresie geometrică}:
\[
  L[i] = \left\lfloor \alpha^{\,i-1} \cdot L_1 + 0{,}5 \right\rfloor, \quad i = 1,\ldots,M-1
\]
unde $\alpha$ este factorul de creștere și $L_1$ este lungimea minimă de istoric
(implicit $\alpha=2$, $L_1=2$, rezultând seria $\{0,2,4,8,16,32,64,128\}$).

Față de o progresie liniară, cea geometrică acoperă eficient un interval larg de scări:
istorii scurte asigură granularitate fină, iar creșterea exponențială permite atingerea
rapidă a dependențelor îndepărtate. Valori extreme ale lui $\alpha$ sunt contraproductive:
$\alpha\geq3$ introduce goluri în acoperire, iar $\alpha\approx1{,}2$ produce tabele
redundante. Factorul $\alpha$ se poate calcula din lungimile extreme ale seriei:
\[
  \alpha = \left(\frac{L(M-1)}{L(1)}\right)^{\!\frac{1}{M-2}}
\]

\subsubsection*{Pragul de actualizare $\theta$}

Predictorul calculează suma algebrică a contoarelor din toate tabelele,
$S = \sum_{i=0}^{M-1} T_i$,
și prezice \textit{taken} dacă $S \geq 0$, altfel \textit{not taken}.
Contoarele sunt actualizate doar la mispredictie sau când $|S| \leq \theta$,
evitând modificări inutile atunci când predictorul este deja sigur.

Valoarea implicită $\theta \approx M$ este justificată de faptul că suma maximă crește
proporțional cu numărul de tabele. Un $\theta$ prea mic conduce la actualizări excesive
(instabilitate, overfitting); un $\theta$ prea mare încetinește adaptarea predictorului.

\subsection{Fluxul de functionare}

\paragraph{Predictie.}
Pentru un PC dat, fiecare tabel $i$ calculeaza un index:
\[
  \mathit{idx}_i = \bigl(PC \oplus \mathit{GHR}[0:L[i]]\bigr) \bmod \mathit{tableSize}
\]
Valoarea contorului de la acel index, $c_i$, este adaugata la suma globala:
\[
  S = \sum_{i=0}^{M-1} c_i[\,\mathit{idx}_i\,]
\]
Decizia finala de predictie este:
\[
  \hat{y} =
  \begin{cases}
    \textit{taken}     & \text{dacă } S \geq 0 \\
    \textit{not taken} & \text{dacă } S < 0
  \end{cases}
\]

\paragraph{Actualizare.}
Dupa rezolvarea branch-ului, predictorul executa trei operatii in ordine:

\begin{enumerate}
    \item \textbf{Actualizarea contoarelor de predictie}: daca a existat
          o mispredictie sau daca predictorul are incredere redusa in
          predictie ($|\mathit{sum}| \leq \theta$), fiecare contor accesat
          este actualizat prin incrementare (pentru branch \textit{taken})
          sau decrementare (pentru branch \textit{not taken}).

          Ideea este ca predictorul invata agresiv doar atunci cand:
          \begin{itemize}
              \item predictia a fost gresita; sau
              \item predictia a fost corecta, dar confidence-ul a fost mic.
          \end{itemize}

          Astfel se evita modificarea inutila a contoarelor atunci cand
          predictorul este deja foarte sigur de decizie.

    \item \textbf{Ajustarea dinamica a lungimii istoriei}
          (\textit{Dynamic History Length Fitting}): mecanismul monitorizeaza
          mispredictiile recente pentru a detecta daca acestea sunt cauzate
          de aliasing (coliziuni intre istorii diferite care acceseaza
          aceeasi intrare din tabel).

          Daca aliasing counter-ul satureaza pozitiv, predictorul considera
          ca istoricul curent este insuficient pentru capturarea corelatiilor
          si activeaza un GHR (\textit{Global History Register}) mai lung
          pentru tabelele cu indici pari. Daca aliasing counter-ul satureaza
          negativ, predictorul revine la un GHR mai scurt pentru reducerea
          zgomotului si a aliasing-ului.

          Comutarea doar pentru jumatate dintre tabele permite coexistenta:
          \begin{itemize}
              \item corelatiilor locale (istorii scurte);
              \item corelatiilor pe termen lung (istorii lungi).
          \end{itemize}

          Astfel predictorul se poate adapta dinamic la diferite aplicatii
          sau chiar la faze diferite ale aceluiasi program.

    \item \textbf{Ajustarea adaptiva a pragului $\theta$}
          (\textit{Adaptive Threshold Fitting}): pragul $\theta$ controleaza
          confidence-ul minim necesar pentru a evita actualizarea predictorului.

          Predictorul produce o suma:

          \[
              \mathit{sum} = \sum_i T_i
          \]

          iar magnitudinea $|\mathit{sum}|$ reprezinta nivelul de incredere
          al predictiei.

          Threshold counter-ul este incrementat la mispredictie si decrementat
          atunci cand predictia este corecta, dar confidence-ul este redus
          ($|\mathit{sum}| \leq \theta$). La saturare, valoarea lui $\theta$
          este ajustata automat:

          \begin{itemize}
              \item cresterea lui $\theta$ face predictorul mai agresiv,
                    permitand mai multe update-uri;
              \item scaderea lui $\theta$ face predictorul mai conservator,
                    reducand update-urile inutile.
          \end{itemize}

          Acest mecanism permite adaptarea dinamica a sensibilitatii
          predictorului la comportamentul curent al aplicatiei si reduce
          atat zgomotul, cat si supra-antrenarea predictorului.
\end{enumerate}

\subsection{Configurarea predictorului}

Implementarea expune trei parametri configurabili, prezentati in
Tabelul~\ref{tab:params}.

\begin{table}[H]
    \centering
    \small
    \begin{tabular}{|l|c|c|c|l|}
        \hline
        \textbf{Parametru} & \textbf{Implicit} & \textbf{Min} & \textbf{Max} & \textbf{Descriere} \\
        \hline
        Numar de tabele ($N$)      & 7    & 4 & 12    & Numarul de tabele de predictie \\
        Marimea tabelelor          & 2048 & 1 & 88000 & Numarul de intrari per tabel \\
        Lungimea contoarelor ($k$) & 4    & 3 & 5     & Precizia contoarelor saturate (biti) \\
        \hline
        $\alpha$ (raport geometric) & 2   & -- & --   & Factorul de crestere al istoriei \\
        $L_1$ (istorie minima)      & 2   & -- & --   & Lungimea istoriei pentru tabelul 1 \\
        $\theta$ initial            & $N$ & -- & --   & Pragul adaptiv de actualizare \\
        \hline
    \end{tabular}
    \caption{Parametrii predictorului OGEHL}
    \label{tab:params}
\end{table}

